\section{Introduction}
\label{sec:intro}

Hardware design remains one of the most error-prone and labor-intensive disciplines in computing. A single bit-width mismatch or missing case in a finite state machine can propagate silently through synthesis, only to surface weeks later as a silicon bug costing millions of dollars to fix. Despite decades of progress in electronic design automation (EDA), the dominant hardware description languages---Verilog and VHDL---still lack the static guarantees that modern software engineers take for granted: dependent types, exhaustive pattern matching, and machine-checked proofs of correctness.

Two largely independent research threads have sought to address this gap. On one hand, \emph{type-safe HDLs} such as Chisel~\citep{bachrach2012chisel}, Bluespec~\citep{nikhil2004bluespec}, and C$\lambda$ash~\citep{baaij2010clash} embed hardware descriptions in host languages with richer type systems, eliminating entire classes of bugs at compile time. On the other hand, \emph{large language models} (LLMs) have demonstrated remarkable ability to generate Verilog from natural language specifications~\citep{thakur2023benchmarking,liu2023verilogeval}, promising to democratize hardware design by lowering the barrier to entry. Yet these two threads remain disconnected: LLM-generated Verilog inherits all the pitfalls of the target language---unintended latches, combinational loops, and width mismatches---while type-safe HDLs still require expert knowledge to use effectively.

We argue that these threads are not merely complementary but \emph{synergistic}: a dependently-typed HDL provides exactly the kind of rich, structured feedback that an LLM agent needs to iteratively converge on correct hardware. In this paper, we present \textsc{Sparkle}, a framework that unifies LLM-driven code generation with formal hardware verification by targeting a hardware DSL embedded in Lean~4~\citep{moura2021lean4}. Given a natural language specification, an LLM coding agent generates hardware in Sparkle's \texttt{Signal} DSL, where Lean's type checker immediately catches width mismatches, the monadic structure prevents combinational loops by construction, and exhaustive pattern matching eliminates unintended latches. When compilation succeeds, the framework automatically extracts synthesizable SystemVerilog, runs functional simulation against reference testbenches, and performs physical synthesis through OpenROAD~\citep{ajayi2019openroad} with the SkyWater 130nm PDK~\citep{edwards2004sky130}. A closed-loop PPA optimization stage then feeds power, performance, and area metrics back to the agent, enabling iterative architectural refinement with formal equivalence guarantees.

The key insight behind \textsc{Sparkle} is that \emph{the compiler is the verifier}. Rather than generating Verilog and hoping it is correct, the agent receives immediate, actionable feedback from Lean's type system at every iteration---feedback that is far more informative than the cryptic warnings of a Verilog linter. This transforms the LLM's generate-and-test loop from a stochastic search over syntactically valid programs into a \emph{type-guided refinement} process that systematically narrows the space of candidate designs.

Our contributions are as follows:
\begin{itemize}[leftmargin=*,itemsep=2pt,topsep=3pt]
    \item We introduce \textsc{Sparkle}, the first framework that combines LLM-driven hardware generation with a dependently-typed HDL, enabling type-guided iterative refinement from natural language to synthesizable silicon.
    \item We design an end-to-end automated pipeline spanning code generation, formal type checking, functional simulation, logic synthesis, and place-and-route, with a closed-loop PPA optimization stage that preserves functional correctness.
    \item We demonstrate that targeting a type-safe HDL substantially improves LLM success rates compared to direct Verilog generation, as compile-time type errors provide structured feedback that guides the agent toward correct designs.
    \item We show that the framework can produce formally verified hardware---including machine-checked equivalence proofs between specifications and optimized implementations---bridging the gap between LLM-generated code and provably correct hardware.
\end{itemize}
